\cvsection{Expériences professionnelles}

\begin{cventries}
  \cventry
    {Stagiaire ingénieur de recherche en IA}
    {INRIA}
    {Bordeaux, France}
    {Mars 2026 -- Septembre 2026}
    {
      \begin{cvitems}
        \item {Utilisation de ressources HPC avec Slurm pour expérimenter sur les LLM, avec installation, configuration et serving de modèles via vLLM sur GPU NVIDIA.}
        \item {Définition d’un protocole expérimental pour comparer Qwen, Gemma et GLM selon la qualité des réponses, la latence, le TTFT et la taille du contexte.}
        \item {Comparaison de plusieurs quantifications pour étudier les compromis entre qualité et performances selon les environnements d’exécution.}
        \item {Analyse des résultats expérimentaux et rédaction en cours d’un article scientifique sur la comparaison des modèles de langage pour notre cas d’usage.}
        \item {Conteneurisation de composants avec Docker, versionnement sous Git et mise en place d’une CI/CD sur GitHub.}
        \item {Évaluation d’OpenClaw dans un environnement Docker isolé et analyse des risques liés à l’accès autonome aux outils et aux fichiers.}
        \item {Gestion d’une application web utilisant des LLM installés localement avec LM Studio sur un Mac Studio, avec installation et configuration des modèles.}
        \item {Conception et développement en autonomie, avec le suivi de mes encadrants, d’un agent IA pour la recherche d’articles scientifiques, à partir des besoins définis avec mon tuteur.}
        \item {Développement de l’agent en Python sans framework d’orchestration, avec appels d’outils, gestion de la mémoire et du contexte, et intégration de l’API OpenAI.}
        \item {Développement de services Python/FastAPI, d’API REST documentées avec OpenAPI et d’une interface React/Next.js, avec persistance MongoDB et pipeline RAG.}
        \item {Automatisation en Python de la collecte de métadonnées d’articles scientifiques en ligne et transmission des résultats de tâches cron par webhooks.}
      \end{cvitems}
    }

  \cventry
    {Stagiaire ingénieur de recherche en apprentissage par renforcement}
    {Université Karunya}
    {Coimbatore, Inde}
    {Juin 2025 -- Sept. 2025}
    {
      \begin{cvitems}
        \item {Développement d’un modèle d’apprentissage par renforcement et de son environnement d’entraînement pour la navigation autonome d’un robot hospitalier.}
        \item {Déploiement et optimisation du logiciel et du modèle sur Raspberry Pi pour des tests en conditions réelles.}
      \end{cvitems}
    }
\end{cventries}
